\begin{vidu}%SBT12.4
Một viên đạn được bắn theo phương nằm ngang từ một khẩu súng đặt ở độ cao 45 m so với mặt đất. Tốc độ của viên đạn khi vừa ra khỏi nòng súng là 250 m/s. Lấy $g=9,8\ m/s^2$.
\begin{itemize}
\item[a)] Sau bao lâu thì viên đạn chạm đất? 
\item[b)] Viên đạn rơi xuống đất cách điểm bắn theo phương nằm ngang bao nhiêu mét?
\item[c)] Ngay trước khi chạm đất, tốc độ của viên đạn bằng bao nhiêu?
\end{itemize}
\haidong{15}
\loigiai{
\begin{itemize}
\item[a)] $t=\sqrt{\dfrac{2h}{g}}=\sqrt{\dfrac{2.45}{9,8}}=3,03\ s$.
\item[b)] $L=v_0t=757,5\ m$.
\item[c)] $v_y=gt=9,8.3,03=29,7\ m/s$.\\
$v=\sqrt{v_x^2+v_y^2}=\sqrt{29,7^2+250^2}\approx 252\ m/s$
\end{itemize}
}
\end{vidu}	
\begin{vidu}%SBT12.5
Một vật được ném theo phương nằm ngang từ độ cao 4,9 m, có tầm xa trên mặt đất $L=5\ m$. Lấy $g=9,8\ m/s^2$.
\begin{itemize}
\item[a)] Tính vận tốc ban đầu.
\item[b)] Xác định tốc độ của vật ngay trước khi chạm đất. 
\item[c)] Xác định hướng chuyển động của viên đạn hợp với phương ngang khi vừa chạm đất.
\end{itemize}
\haidong{15}
\loigiai{
\begin{itemize}
\item[a)] $L=v_0\sqrt{\dfrac{2h}{g}}\Rightarrow v_0=L.\sqrt{\dfrac{g}{2h}}=5.\sqrt{\dfrac{9,8}{2.4,9}}=5\ m/s$.
\item[b)] \begin{align*}
&v_x=v_0=5\ m/s\\
&v_y=gt=\sqrt{2gh}=9,8\ m/s\\
&v=\sqrt{v_x^2+v_y^2}=\approx 11\ m/s 
\end{align*}
\item[c)] $\tan\alpha=\dfrac{v_y}{v_x}=1,96\Rightarrow \alpha \approx 63^\circ$ 
\end{itemize}
}
\end{vidu}	

\begin{vidu} %[1P2K9-2][Loai1]
Một người ném quả bóng theo phương nằm ngang với vận tốc đầu $v_0 = 15$ m/s và rơi xuống cách mình một khoảng 45 m theo phương ngang. Thời gian quả bóng rơi là
\choice
{5 s}
{4 s}
{\True 3 s}
{2 s}
\haidong{4}
\loigiai{ 
Ta có thời gian rơi cần tính: $x = v_0t \Rightarrow t = \dfrac{x}{v_0} = \dfrac{45}{15} = 3$ s.
}
\end{vidu}
\begin{vidu} %[1P2-9.2K][Loai1]
Một hòn bi lăn dọc theo một cạnh của một mặt bàn hình chữ nhật nằm ngang cao $h = 1,25$ m. Khi ra khỏi mép bàn, nó chạm nền nhà tại điểm cách mép bàn $L = 1,50$ m theo phương ngang? Bỏ qua sức cản của không khí và lấy $g = 10\ m/s^2$. Thời gian từ lúc rời mép bàn đến khi chạm đất của hòn bi là
\choice
{0,35 s}
{0,125 s}
{\True 0,5 s}
{0,25 s}
\haidong{5}
\loigiai{      
\begin{itemize}
\item Thời gian rơi $t = \sqrt {\dfrac{2h}{g}} = \sqrt {\dfrac{{2.1,25}}{10}} = 0,5\ s.$
\end{itemize}}
\end{vidu}
\begin{vidu} %[1P2-9.2K][Loai1]
Một quả bóng được ném theo phương nằm ngang từ đỉnh một vách đá trên mặt đất với vận tốc 40 m/s. Bỏ qua sức cản của không khí và lấy $g = 10\ m/s^2$. Sau 3 s kể từ lúc ném, vận tốc của quả bóng đó bằng
\choice
{30 m/s}
{40 m/s}
{\True 50 m/s}
{70 m/s}
\haidong{5}
\loigiai{      
\begin{itemize}
\item Đây là chuyển động ném ngang: $v_x = v_0 = 40$ m/s; $v_y = gt = 10.3 = 30$ m/s.
\item Do đó: $v_t = \sqrt{v_x^2 + v_y^2} = 50\ m/s.$
\end{itemize}}
\end{vidu}
\begin{vidu} %[1P2K9-2][Loai1]
Một máy bay bay theo phương ngang ở độ cao 10 km với tốc độ 720 km/h.  Bỏ qua sức cản của không khí và lấy $g = 10\ m/s^2$. Viên phi công phải thả quả bom từ xa, cách mục tiêu theo phương ngang bao nhiêu để quả bom rơi trúng mục tiêu?
\choice
{\True 8,94 km}
{6,32 km}
{89,44 km}
{4,47 km}
\haidong{5}
\loigiai{ 
\begin{itemize}
\item Để quả bom rơi trúng mục tiêu thì viên phi công phải thả quả bom từ cách xa mục tiêu một đoạn bằng tầm ném xa của quả bom: 
$L=x_{\max}=v_0\sqrt{\dfrac{2h}{g}}\approx 8,94\ km$
\item Chú ý: đổi đơn vị $720\ km/h = 200\ m/s;\ 10\ km = 1000\ m$.
\end{itemize}}
\end{vidu}
\begin{vidu}
Từ đỉnh một ngọn tháp cao $80\ m$, một quả cầu được ném theo phương ngang với vận tốc đầu $20\ m/s$. Bỏ qua lực cản của không khí và lấy $g = 10\ m/s^2$. Độ dịch chuyển thành phần của quả cầu là 
\choice
{\True $d_y=5 t^2;\ d_x=20 t$}
{$d_y=10 t^2;\ d_x=20 t$}
{$d_y=5 t^2;\ d_x=40 t$}
{$d_y=10 t^2;\ d_x=10 t$}
\haidong{7}
\loigiai{
\immini[thm]{
Độ dịch chuyển của quả cầu
$
\left\{\begin{array} {l} 
{d_x = v _ {0} t} \\
{d_y = \frac {1} {2} g t ^ {2}}
\end{array}\right.  \Rightarrow \left\{\begin{array}{l}
d_x=20 t \\
d_y=5 t^2
\end{array}\right.
$}{\includegraphics[width=5cm]{images/10_2}}
}
\end{vidu}
%\loai{Viết phương trình quỹ đạo}
\begin{vidu} %[1P2-9.2K][Loai1]
Một máy bay bay với vận tốc không đổi $v_0 = 120$ m/s theo phương nằm ngang ở độ cao h so với mặt đất và thả một vật. Bỏ qua sức cản của không khí và lấy $g = 10\ m/s^2$. Phương trình quỹ đạo của vật khi chọn gốc tọa độ O ở điểm thả vật, Ox hướng theo $\vec v_{0}$; Oy hướng thẳng đứng xuống dưới là
\choice
{\True $y = \dfrac{1}{2880}x^2$}
{$y = \dfrac{1}{1440}x^2$}
{ $y = \dfrac{1}{120}x^2$}
{ $y = \dfrac{1}{240}x^2$}
\haidong{8}
\loigiai{      
Ta có phương trình quỹ đạo: $y = \dfrac{1}{2}gt^2 = \dfrac{1}{2}g\left( {\dfrac{x}{{v_0}}}\right)^2 = \dfrac{1}{2880}x^2.$
}
\end{vidu}
